\begin{textsample}{NEG Dim 1 – al – Score 4.00 – al/al\_inf\_22.txt}  \label{ex:f1_neg_135}
6.
AVALIAÇÃO E REGISTRO Conforme é previsto nas DCNEI, “ as instituições de Educação Infantil devem criar procedimentos para acompanhamento do trabalho pedagógico e para avaliação do desenvolvimento das crianças, sem objetivo de seleção, promoção ou classificação ” ( BRASIL, 2010, p. 29 ).
Está claro que esta avaliação não deve ser para promover ou classificar, mas garantir o acompanhamento do desenvolvimento da criança, conforme supracitado neste documento.
É necessário que as IEI criem mecanismos e procedimentos de acompanhamento por meio de : - Observação crítica e criativa das atividades, das brincadeiras e interações das crianças no cotidiano ; - Utilização de múltiplos registros realizados por adultos e crianças ( relatórios, fotografias, desenhos, álbuns etc. ) ; - Documentação específica que permita às famílias conhecer o trabalho da instituição junto às crianças e os processos de desenvolvimento e aprendizagem da criança na Educação Infantil.
Ao elaborar esta documentação, deve se considerar os campos de experiências e os direitos de aprendizagem defendidos na BNCC, de modo que estes sejam contemplados no planejamento docente e avaliados durante a prática pedagógica e seus efeitos junto às crianças.
Importante reiterar que a avaliação mediante acompanhamento e registro do desenvolvimento das crianças não deve ter o objetivo de promoção e menos ainda de retenção.
Portanto, se faz necessária planejar o processo de transição entre estas duas etapas ( EI e EF ), considerando a progressão da aprendizagem das crianças.
Realizada com o propósito de informar o professor e a criança sobre a progressão da aprendizagem, possibilitando reformulações e assegurando o alcance dos objetivos, a avaliação na Educação Infantil demonstra como as crianças se desenvolvem.
Por isso, deve ser realizada de maneira contínua, no dia-a-dia, e também em oportunidades regulares, incluindo o uso de instrumentos mais formais ( relatórios, fichas avaliativas e portfólios ) pautados na observação, no registro diário e na singularidade de cada criança.
Na observação, deve analisar as características individuais das crianças, como elas lidam com as outras, como elas participam das experiências de aprendizagem, suas conquistas, dificuldades e possibilidades de desenvolvimento.
O registro deve ser um momento de anotações de tudo que foi observado, suas falas importantes, suas expressões, a maneira que se processa a construção do pensamento do conhecimento, sua maneira de interagir com as experiências, quais os direitos de aprendizagem que são garantidos.
O portfólio, quando bem construído, reproduz a realidade da turma e de cada aluno.
Ele pode conter relatórios, fotos, atividades das crianças, fichas descritivas, vídeos com experiências das brincadeiras e interações.
O importante do portfólio que pode ser socializado a qualquer momento com a comunidade escolar e pode servir como objeto de estudo para o professor na construção de planejamentos posteriores.
Acompanhar o monitoramento do desenvolvimento da criança na Educação Infantil exige progressão e articulação de ações diariamente.
O professor na condição de avaliador propõe, observa, introduz e redireciona seu planejamento no intuito de : - Estimular a criatividade e curiosidade da criança nas atividades diárias da instituição ; - Valorizar a diversidade, interesses e possibilidades dos direitos de aprendizagens : conviver, explorar, brincar, participar, expressar -se e conhecer -se no seu mundo. - Medir o pensamento da criança com possibilidades e brincadeira desafiadoras e possíveis de serem executadas. - Fortalecer o reconhecimento das diferenças e das diversidades étnicas, culturais, política, ambiental e regional do povo alagoano. - Registrar os aspectos individuais do desenvolvimento das crianças em todos seus aspectos : cognitivos, social, ético, estético, físico e psicológico.

% matched lemmas: 
\end{textsample}
